\section{Introduction}
In this chapter, GPAT is used to conduct a series of studies to explore the potential for climate benefits through controlled formation flight of commercial aircraft, and the resulting overlap of their exhaust plumes. As explained in detail in section \ref{Saturation}, microphysical and gas-phase photochemical saturation effects can occur under high emissions concentrations in the Upper Troposphere (UT), where aircraft typically cruise. This means that inside a fresh exhaust plume, or in a region where many aircraft plumes are clustered, nonlinear plume-scale saturation effects can occur, influencing the fate of key chemical species and aircraft contrails. 

Aircraft nitrogen oxide (\ce{NOx} = NO + \ce{NO2}) emissions are responsible for around 16~\% of aviation climate impact, primarily due to the short-term production of ozone, a potent greenhouse gas. Emissions of water vapour and particulates from the aircraft exhaust can generate condensation trails (contrails), which are thought to have a net warming effect that is responsible for around 50~\% of aviation-induced climate change. There are two primary saturation effects that occur in high-emissions regions at a particular flight level in the UT, which can reduce the climate impact of aviation \ce{NOx} and contrails; when critical \ce{NOx} threshold concentrations are surpassed, ozone production efficiency decreases with increasing \ce{NOx}. It is thought this could reduce the net-\ce{NOx}  Furthermore, when persistent contrails are formed in a close proximity to one another, they experience mutually inhibited growth due to competing for finite ambient water vapour. In both cases, this means that flying aircraft close together and overlapping their exhaust plumes can provide climate benefits over the alternative scenario, where plumes evolve independently of each other.

The following series of studies presented in this chapter focus on \ce{NOx} saturation, building on many years of experience and previous studies published by the University of Bristol ACRG (e.g. \cite{Derwent2024}). It is also a direct continuation of the research presented thus far in this thesis, and serves as a prime use case for the GPAT modelling framework.

\subsection{\ce{NOx}-Saturated Chemistry Primer}
When aircraft plumes overlap, \ce{NOx} emissions accumulate, and can tip the local domain from a \ce{NOx}-limited to a \ce{NOx}-saturated (or VOC-limited) regime. In the \ce{NOx}-limited regime, hydroxyl radicals (OH) typically react with volatile organic compounds (VOCs), driving the daytime odd hydrogen cycle (\ce{HOx} = OH + \ce{HO2}). This process generates organic peroxy radicals (\ce{RO2}), which contribute to ozone formation through the following mechanism:

\reaction[RH+OH]{RH + OH -> R + H2O}
\reaction[R+O2]{R + O2 -> RO2}

Aircraft \ce{NOx} emissions are composed mostly of nitric oxide (\ce{NO} $\approx$ 95~\%), with the remainder being nitrogen dioxide (\ce{NO2} $\approx$ 5~\%) \cite{Klemm1998}. Aircraft NO speeds up the catalytic oxidation of VOCs and reduces ozone destruction by \ce{HOx} through reaction with peroxy radicals to form \ce{NO2}:

\reaction[HO2+NO]{HO2 + NO -> OH + NO2}
\reaction[RO2+NO]{RO2 + NO -> RO + NO2}

NO to \ce{NO2} conversion is followed by photolysis of \ce{NO2}, which results in the production of atomic oxygen, and subsequently ozone \cite{Fittschen2019} through \eqref{NO2+hv} and \eqref{O3P+O2+M}.

\reaction[NO2+hv]{NO_2 + h\nu (\lambda $<$ 420 nm) -> NO + O(^3P)}
\reaction[O3P+O2+M]{O(^3P) + O_2 + M -> O_3 + M}

When \ce{NO2} concentrations exceed a threshold level, the local system transitions from \ce{NOx}-sensitive to \ce{NOx}-saturated conditions due to a high \ce{NOx}-VOC ratio meaning that termination reactions of NO and \ce{NO2} to reservoir species (e.g. \ce{HONO}, \ce{HNO3}, \ce{HO2NO2}) are often favoured over catalytic ozone formation. These reactions act as a net sink for \ce{NOx} due to low reactivity and eventual loss through deposition. NO titration then becomes a dominant pathway for ozone destruction, reducing the net \ce{O3} production efficiency (see section \ref{Saturation}: \ce{NOx}-Saturated Regime for more on this). 

%Talk about Nussbaumer and potential contradictions

\subsection{Exploiting \ce{NOx} Saturation through Plume Overlap}
The dynamical evolution of aircraft exhaust plumes is a phenomenon often overlooked in global and regional studies of the atmospheric effects resulting from aviation emissions. The dispersion and advection of the exhaust mass happens over timescales of hours, meaning emissions concentrations can remain much higher than background inside the plume during this time. Global and regional models assume instantaneous dispersion (ID), in which dynamics are reduced to homogeneous distribution across the occupied grid cell at the time of emission \cite{Paoli2011}. This assumption precludes any attempt to model saturation processes which occur on the subgrid scale. Such effects are exacerbated when the aircraft transit frequency through a region is high. The following studies begin to address the nuance in the ozone mitigation potential of formation flight due to plume intersections. When many aircraft plumes overlap, the \ce{NOx} emissions in a localised region of the atmosphere can reach saturated levels, but this highly depends on a number of factors such as time, location, number of aircraft, separation minima, and metrics used.

Firstly, five representative case scenarios are presented, which show \ce{NOx} saturation potential for high traffic density hotspot locations. A Background Study is then carried out, which tests a range of metrics indicative of ozone sensitivity across all locations and months of the year. The ID Vs Plume Study addresses the difference in outputs between an instant dilution modelling approach, and the GPAT plume aggregation approach, to examine the differences in ozone output for a range of aircraft configurations. In the Sensitivity Study, five case scenarios are tested under a single aircraft Vs. five aircraft configuration to explore the region's response to \ce{NOx} saturation. Finally, a Parametric Study is conducted to investigate the variation in saturation benefits under a range of unique aircraft configurations (no. of aircraft and separation distances).